%-------------------------
% Hoja de Vida, una plantilla de currículum en latex para desarrolladores
% Autor : Christian Bueno
% Basado en: https://github.com/christianbueno1/latex-pilot
% Licencia : MIT
%-------------------------

%------------PAQUETES----------------
\documentclass[a4paper,11pt]{article}

% Paquetes de diseño y tipografía
\usepackage{titlesec}     % formato personalizado de secciones
\usepackage{enumitem}     % control sobre listas (itemize, enumerate)
\usepackage{needspace}    % evita cortes de página incómodos
\usepackage{fancyhdr}     % personalización de encabezados y pies de página
\usepackage[empty]{fullpage} % márgenes de página personalizados
\usepackage{array}        % columnas personalizadas en tablas
\usepackage{booktabs}     % líneas de tabla más sutiles y legibles

% Enlaces y formato de texto
\usepackage[hidelinks]{hyperref} % enlaces clicables sin bordes
\usepackage[normalem]{ulem}      % subrayado con soporte para saltos de línea

% Soporte de color (versión moderna)
\usepackage{xcolor}       % gestión avanzada de color

% Soporte de idioma
\usepackage[spanish,es-noquoting]{babel} % reglas tipográficas del español
%-----------------------------------------

\input glyphtounicode % convierte nombres de glifos a unicode
\pdfgentounicode=1 % asegura que los PDFs generados sean legibles por ATS

%----------OPCIONES DE FUENTE-------------------
\usepackage[default]{sourcesans} % usa la fuente Source Sans
\urlstyle{same} % los URLs usan la misma fuente que el documento
%-----------------------------------------

%----------OPCIONES DE MÁRGENES-----------------
\pagestyle{fancy} % estilo de página personalizable con fancyhdr
\fancyhf{} % limpia todos los campos de encabezado y pie de página

\renewcommand{\headrulewidth}{0in} % grosor de línea bajo el encabezado a cero
\renewcommand{\footrulewidth}{0in} % grosor de línea sobre el pie de página a cero

\setlength{\tabcolsep}{0in} % grosor del separador de columnas en tablas a cero
\setlength{\footskip}{4.08003pt} % espacio mínimo para pie de página (evita warnings)

% origen del documento es una pulgada desde arriba y desde la izquierda
% oddsidemargin y evensidemargin se refieren al margen izquierdo
% el margen derecho se ajusta indirectamente usando oddsidemargin
\addtolength{\oddsidemargin}{-0.5in}
\addtolength{\topmargin}{-0.5in}

\addtolength{\textwidth}{1.0in} % aumenta el ancho del área de texto en una pulgada
\addtolength{\textheight}{1.0in} % aumenta la altura del área de texto en una pulgada

\raggedbottom{} % todas las páginas tienen la altura de la página actual, sin espacio vertical extra
\raggedright{} % todas las páginas tienen el ancho de la página actual, sin espacio horizontal extra
%------------------------------------------

%--------COMANDOS DE SECCIÓN---------
\titleformat{\section}
  {\scshape\large}{}
    {0em}{\color{blue}}[\color{black}\titlerule\vspace{0pt}]
%-------------------------------------

%--------REDEFINICIONES----------------
\renewcommand{\ULdepth}{2pt}
%-------------------------------------

%--------COMANDOS PERSONALIZADOS--------------
\newcommand{\resumeItem}[1]{
  \item\small{#1}
}

% rightmargin ajustado a 0.11in para evitar desbordamiento de resumeItem
\newcommand{\resumeItemListStart}{\begin{itemize}[rightmargin=0.11in]}
\newcommand{\resumeItemListEnd}{\end{itemize}}

% resumeQuadHeading renderiza cuatro elementos en dos columnas
\newcommand{\resumeQuadHeading}[4]{
  \item
  \begin{tabular*}{0.96\textwidth}[t]{l@{\extracolsep{\fill}}r}
    \textbf{#1} & #2 \\
    \textit{\small#3} & \textit{\small #4} \\
  \end{tabular*}
}

% lefmargin para indentación izquierda de 0.15in para resumeItems
\newcommand{\resumeHeadingListStart}{
  \begin{itemize}[leftmargin=0.15in, label={}]
}
\newcommand{\resumeHeadingListEnd}{\end{itemize}}

% comandos para detalles de contacto
\newcommand{\name}[1]{\textbf{\Huge #1}}
\newcommand{\socialLinks}[1]{#1}
\newcommand{\separator}{$|$}
%-------------------------------------------
